\section{Results}

\subsection*{Empirical objects}

The inferential object is the utility curve $U(\tau)$, defined by linear access
to future-relevant targets from present representations across horizons
$\tau$. The summary observables are inferential half-life, utility area, and
the adaptation burden target
\[
  R^*(m,\pi) = \min_w R(m,\pi,w),
\]
the minimum post-drift loss area over adaptation-window choices for protocol
$\pi$. This definition lets the burden comparison absorb the non-monotone window
effects already present in recurrent multiscale adaptation.

\subsection*{Family-level regime separation}

The first empirical regularity is a family-level ordering of burden regimes.
Under the fast four-seed benchmark with the full window grid
$w \in \{1,16,64\}$, the abrupt-drift mean burdens are approximately
$2.1, 2.2, 3.2,$ and $12.2$ for feedforward local, feedforward multiscale,
recurrent local, and recurrent multiscale variants respectively. Under gradual
drift the means are approximately $2.8, 2.7, 9.3,$ and $22.3$ in the same
order. The qualitative hierarchy therefore persists across the two tested drift
types: feedforward variants occupy the low-burden regime, recurrent local lies
intermediate, and recurrent multiscale is the highest-burden regime.

Effective spectral rank is the strongest pooled regime separator in the current
dataset. On the pooled fast benchmark ($n=32$ cases),
$\mathrm{Spearman}(\mathrm{rank},R^*) \approx 0.74$, whereas
$\mathrm{Spearman}(\Delta_v,R^*) \approx 0.02$. A pooled additive linear model
using z-scored rank and $\Delta_v$ reaches $R^2 \approx 0.38$, and adding the
interaction term rank$\times\Delta_v$ improves fit only marginally
($\Delta R^2 \approx 0.005$). At the pooled level, effective rank therefore
functions as a family-level regime coordinate, while $\Delta_v$ does not act as
a universal burden predictor.

\subsection*{Within-family scale structure in recurrent local models}

The second empirical regularity appears inside a fixed family. In the recurrent
local variant, a targeted capacity sweep over
$\mathrm{hidden} \in \{32,48,64,96\}$ and
$\mathrm{depth} \in \{2,3,4\}$ under abrupt drift yields
$\mathrm{Spearman}(\Delta_v,R^*) \approx 0.71$, while
$\mathrm{Spearman}(\mathrm{rank},R^*) \approx 0.36$. This is not a
low-variance artifact: effective rank still spans approximately $17.6$ to
$65.5$ across the same sweep, yet its Pearson correlation with $R^*$ is only
moderate ($\approx 0.43$). Within this recurrent local regime, dynamic
reactivity is therefore the stronger burden coordinate.

The distinction between those two levels is the main empirical result. Rank
varies enough within the recurrent local family to rule out a trivial variance
explanation, but it still does not dominate burden there. The data support a
hierarchical structure: family choice sets the degradation regime, while
within-family scale variation shifts burden through the activation-velocity
channel.

\subsection*{Cross-drift prognosis remains sample-limited}

The prognosis layer is feasible but not yet closed. On the fast four-seed run, a
gradual-to-abrupt ridge forecaster based on routing alignment, deep incremental
utility area, inferential half-life difference, and peak-gap information reaches
$\rho \approx 0.66$ but negative $R^2$ and does not dominate the strongest
single-feature baseline. Effective rank alone gives the best rank correlation in
that setting ($\rho \approx 0.89$), while $\Delta_v$ yields weaker pooled
generalization. The current evidence therefore supports prognosis as a viable
layer, but not yet as a closed predictive theorem for burden across drift types.

\subsection*{Architecture-moderated rank effects remain open}

Exploratory diagnostics suggest that the semantics of spectral rank may depend on
architecture. In reduced-budget within-variant tests, the recurrent local slice
shows a strong positive relation between $\Delta_v$ and burden, whereas
feedforward local behavior does not follow the same pattern and can even display
sign reversals for rank in very small samples. A pooled moderation fit with an
architecture indicator and rank$\times$architecture interaction improves fit
modestly in a reduced dataset, but the coefficients remain unstable at that
sample size. The supported conclusion is therefore hierarchical rather than fully
mechanistic: rank separates families, and $\Delta_v$ controls burden inside the
recurrent local regime.
